\section{Method}

Keep paper-specific definitions in \texttt{main.tex}. For example,
\texttt{\string\method} is defined by the example paper rather than by the
class, because every project will use different method names and symbols.

\subsection{Writing Modules}

Use the section files in \texttt{secs/} to keep long manuscripts readable.
Small papers can also place all content directly in \texttt{main.tex}. The class
does not require a particular project structure.

\subsection{Math and References}

Common math packages are preloaded. For instance, a method can optimize
\begin{equation}
  \mathcal{L} = \lambda_1 \mathcal{L}_{\mathrm{1}}
    + \lambda_2 \mathcal{L}_{\mathrm{2}},
  \label{eq:loss}
\end{equation}
where \(\lambda_1\) and \(\lambda_2\) balance the two loss terms.
\eqnref{eq:loss} is referenced with \texttt{\string\eqnref}; citations use
\texttt{natbib}, as in \cite{lamport1994latex,mittelbach2004latex}.

\subsection{Algorithm Modules}

Algorithm blocks can live in \texttt{algs/} and be included where they are
discussed. \algref{alg:example} shows the expected structure.

\begin{algorithm}[t]
  \caption{Example Training Loop}
  \label{alg:example}
  \begin{algorithmic}[1]
    \Require Dataset \(\mathcal{D}\), model \(f_\theta\), learning rate \(\eta\)
    \Ensure Updated parameters \(\theta\)
    \For{each mini-batch \(B \subset \mathcal{D}\)}
      \State Compute predictions \(\hat{y} \gets f_\theta(B)\)
      \State Evaluate loss \(\mathcal{L}(\hat{y}, y)\)
      \State Update parameters \(\theta \gets \theta - \eta \nabla_\theta \mathcal{L}\)
    \EndFor
  \end{algorithmic}
\end{algorithm}



\subsection{Implementation Details}

Describe the concrete choices that make the method reproducible: network
architecture, optimizer and schedule, batch size, number of training steps,
and any preprocessing applied to the input data. Readers should be able to
reimplement the method from this subsection alone, without consulting the
released code. Where a choice was tuned rather than fixed in advance, state
the search range and the criterion used to select the final value.
